% GENERATED FILE. Edit canonical lessons, then run npm run generate:books.
% canonical-source-hash: fnv1a64:4009d909fbe5f0b3
% canonical-lessons: JA-C279-kazan, JA-C279-funka, JA-C279-saigai, JA-C279-fusegu, JA-C279-oukii, JA-R279-review-274-279

\chapter{The Mountain Erupts}
\label{ch:ja-the-mountain-erupts}
\hlchaptermodality{279}

\begin{chapteropening}
\textbf{By the end of this chapter:} \emph{I can say what a volcano did and how we kept safe.}

Complete the last lesson of chapter 279: I can say what a volcano did and how we kept safe.
\end{chapteropening}

\section[\texorpdfstring{\ja{かざん} (kazan)}{かざん (kazan)}]{\ja{かざん} (kazan) — a volcano}

\label{lesson:JA-C279-kazan}



\begin{quote}
Before the new one: say the Japanese for a tornado, then the Japanese for to lean.
\end{quote}

\subsection*{You'll want to know: \ja{かざん}}
\textbf{\ja{かざん}} — \emph{kazan} — \textquotedblleft{}a volcano\textquotedblright{}.

\begin{quote}
\textbf{\ja{日本}\ja{には} \ja{かざんが} \ja{おおいです。}} — \emph{nihon ni wa kazan ga ōi desu} — \textquotedblleft{}Japan has many volcanoes.\textquotedblright{}
\end{quote}

\begin{cousinweb}
\textbf{\ja{かざん}} (\emph{kazan}) is Sino-Japanese (kango): two Chinese characters read with their on'yomi, readings borrowed from Chinese. They are \emph{ka}, \textquotedblleft{}fire\textquotedblright{} (Middle Chinese \emph{xwaX}) and \emph{zan}, \textquotedblleft{}mountain\textquotedblright{} (\emph{srean}) (KANJIDIC2; Middle Chinese in Baxter's transcription). The second is voiced from \emph{san} to \emph{zan}.
\end{cousinweb}

\begin{grammarlens}[title={What you've built}]
One more word for telling what happened.
\end{grammarlens}

\subsection*{Your turn}
\begin{itemize}
\raggedright
  \item \emph{Say it:} \emph{kazan}
  \item \emph{Say it:} \emph{kazan}, once more
  \item \emph{Say it:} \emph{nihon ni wa kazan ga ōi desu} — Japan has many volcanoes.
  \item \emph{Recall:} say the Japanese for a tornado, then the Japanese for to lean, then say \emph{kazan} again
\end{itemize}

\begin{tcolorbox}[breakable,colback=teal!4,colframe=teal!35!black,title={Before you move on}]
What does \ja{かざん} mean? (\textquotedblleft{}A volcano\textquotedblright{}.) Say it once more.
\end{tcolorbox}
\section[\texorpdfstring{\ja{ふんか} (funka)}{ふんか (funka)}]{\ja{ふんか} (funka) — an eruption}

\label{lesson:JA-C279-funka}



\begin{quote}
Before the new one: say the Japanese for to lean, then the Japanese for a volcano.
\end{quote}

\subsection*{You'll want to know: \ja{ふんか}}
\textbf{\ja{ふんか}} — \emph{funka} — \textquotedblleft{}an eruption\textquotedblright{}.

\begin{quote}
\textbf{\ja{その} \ja{かざんが} \ja{ふんかしました。}} — \emph{sono kazan ga funka shimashita} — \textquotedblleft{}That volcano erupted.\textquotedblright{}
\end{quote}

\begin{cousinweb}
\textbf{\ja{ふんか}} (\emph{funka}) is Sino-Japanese (kango): two Chinese characters read with their on'yomi, readings borrowed from Chinese. They are \emph{fun}, \textquotedblleft{}erupt\textquotedblright{} (Middle Chinese \emph{phwonH}) and \emph{ka}, \textquotedblleft{}fire\textquotedblright{} (\emph{xwaX}) (KANJIDIC2; Middle Chinese in Baxter's transcription).
\end{cousinweb}

\begin{grammarlens}[title={What you've built}]
One more word for telling what happened.
\end{grammarlens}

\subsection*{Your turn}
\begin{itemize}
\raggedright
  \item \emph{Say it:} \emph{funka}
  \item \emph{Say it:} \emph{funka}, once more
  \item \emph{Say it:} \emph{sono kazan ga funka shimashita} — That volcano erupted.
  \item \emph{Recall:} say the Japanese for to lean, then the Japanese for a volcano, then say \emph{funka} again
\end{itemize}

\begin{tcolorbox}[breakable,colback=teal!4,colframe=teal!35!black,title={Before you move on}]
What does \ja{ふんか} mean? (\textquotedblleft{}An eruption\textquotedblright{}.) Say it once more.
\end{tcolorbox}
\section[\texorpdfstring{\ja{さいがい} (saigai)}{さいがい (saigai)}]{\ja{さいがい} (saigai) — a disaster; a calamity}

\label{lesson:JA-C279-saigai}



\begin{quote}
Before the new one: say the Japanese for a volcano, then the Japanese for an eruption.
\end{quote}

\subsection*{You'll want to know: \ja{さいがい}}
\textbf{\ja{さいがい}} — \emph{saigai} — \textquotedblleft{}a disaster; a calamity\textquotedblright{}.

\begin{quote}
\textbf{\ja{さいがいの} \ja{ばあいは、がっこうに} \ja{ひなんします。}} — \emph{saigai no baai wa, gakkō ni hinan shimasu} — \textquotedblleft{}In a disaster, we take shelter at the school.\textquotedblright{}
\end{quote}

\begin{cousinweb}
\textbf{\ja{さいがい}} (\emph{saigai}) is Sino-Japanese (kango): two Chinese characters read with their on'yomi, readings borrowed from Chinese. They are \emph{sai}, \textquotedblleft{}disaster\textquotedblright{} (Middle Chinese \emph{tsoj}) and \emph{gai}, \textquotedblleft{}harm\textquotedblright{} (\emph{hajH}) (KANJIDIC2; Middle Chinese in Baxter's transcription).
\end{cousinweb}

\begin{grammarlens}[title={What you've built}]
One more word for telling what happened.
\end{grammarlens}

\subsection*{Your turn}
\begin{itemize}
\raggedright
  \item \emph{Say it:} \emph{saigai}
  \item \emph{Say it:} \emph{saigai}, once more
  \item \emph{Say it:} \emph{saigai no baai wa, gakkō ni hinan shimasu} — In a disaster, we take shelter at the school.
  \item \emph{Recall:} say the Japanese for a volcano, then the Japanese for an eruption, then say \emph{saigai} again
\end{itemize}

\begin{tcolorbox}[breakable,colback=teal!4,colframe=teal!35!black,title={Before you move on}]
What does \ja{さいがい} mean? (\textquotedblleft{}A disaster; a calamity\textquotedblright{}.) Say it once more.
\end{tcolorbox}
\section[\texorpdfstring{\ja{ふせぐ} (fusegu)}{ふせぐ (fusegu)}]{\ja{ふせぐ} (fusegu) — to prevent; to keep off}

\label{lesson:JA-C279-fusegu}



\begin{quote}
Before the new one: say the Japanese for an eruption, then the Japanese for a disaster.
\end{quote}

\subsection*{You'll want to know: \ja{ふせぐ}}
\textbf{\ja{ふせぐ}} — \emph{fusegu} — \textquotedblleft{}to prevent; to keep off\textquotedblright{}.

\begin{quote}
\textbf{\ja{みんなで} \ja{じこを} \ja{ふせいで} \ja{ください。}} — \emph{minna de jiko o fuseide kudasai} — \textquotedblleft{}Please all help to prevent accidents.\textquotedblright{}
\end{quote}

Polite \textbf{\ja{ふせぎます}}, \emph{fusegimasu}; a godan verb. Past (polite): \textbf{\ja{かさが} \ja{あめを} \ja{ふせぎました。}}, \emph{kasa ga ame o fusegimashita}.

\begin{cousinweb}
\textbf{\ja{ふせぐ}} (\emph{fusegu}) is a native Japanese word (wago), inherited rather than borrowed from Chinese; its kanji was attached later as a kun'yomi (JMdict).
\end{cousinweb}

\begin{grammarlens}[title={What you've built}]
One more word for telling what happened.
\end{grammarlens}

\subsection*{Your turn}
\begin{itemize}
\raggedright
  \item \emph{Say it:} \emph{fusegu}
  \item \emph{Say it:} \emph{fusegu}, once more
  \item \emph{Say it:} \emph{minna de jiko o fuseide kudasai} — Please all help to prevent accidents.
  \item \emph{Recall:} say the Japanese for an eruption, then the Japanese for a disaster, then say \emph{fusegu} again
\end{itemize}

\begin{tcolorbox}[breakable,colback=teal!4,colframe=teal!35!black,title={Before you move on}]
What does \ja{ふせぐ} mean? (\textquotedblleft{}To prevent; to keep off\textquotedblright{}.) Say it once more.
\end{tcolorbox}
\section[\texorpdfstring{\ja{おおきい} (ōkii)}{おおきい (ōkii)}]{\ja{おおきい} (ōkii) — big; large}

\label{lesson:JA-C279-oukii}



\begin{quote}
Before the new one: say the Japanese for a disaster, then the Japanese for to prevent.
\end{quote}

\subsection*{You'll want to know: \ja{おおきい}}
\textbf{\ja{おおきい}} — \emph{ōkii} — \textquotedblleft{}big; large\textquotedblright{}.

\begin{quote}
\textbf{\ja{おおきい} \ja{じしんが} \ja{ありました。}} — \emph{ōkii jishin ga arimashita} — \textquotedblleft{}There was a big earthquake.\textquotedblright{}
\end{quote}

An i-adjective.

\begin{cousinweb}
\textbf{\ja{おおきい}} (\emph{ōkii}) is a native Japanese word (wago), inherited rather than borrowed from Chinese; its kanji was attached later as a kun'yomi (JMdict).
\end{cousinweb}

\begin{grammarlens}[title={What you've built}]
One more word for telling what happened.
\end{grammarlens}

\subsection*{Your turn}
\begin{itemize}
\raggedright
  \item \emph{Say it:} \emph{ōkii}
  \item \emph{Say it:} \emph{ōkii}, once more
  \item \emph{Say it:} \emph{ōkii jishin ga arimashita} — There was a big earthquake.
  \item \emph{Recall:} say the Japanese for a disaster, then the Japanese for to prevent, then say \emph{ōkii} again
\end{itemize}

\begin{tcolorbox}[breakable,colback=teal!4,colframe=teal!35!black,title={Before you move on}]
What does \ja{おおきい} mean? (\textquotedblleft{}Big; large\textquotedblright{}.) Say it once more.
\end{tcolorbox}
\section[(dialogue)]{Looking back: from When the Earth Shook to The Mountain Erupts}

\label{lesson:JA-R279-review-274-279}



\begin{quote}
Say the Japanese for an earthquake and to crumble.
\end{quote}

\subsection*{Your turn}
Say these aloud:

\begin{itemize}
\raggedright
  \item an earthquake, to crumble, to roll, frightening, evacuation
  \item a tsunami, to flow, to wash away, to float, once more
  \item heavy rain, a flood, to overflow, a cliff, awful
  \item heavy snow, to accumulate, a snowstorm, a power cut, to go numb with cold
  \item a flash of lightning, to strike, violent, a tornado, to lean
  \item a volcano, an eruption, a disaster, to prevent, big
\end{itemize}

\begin{tcolorbox}[breakable,colback=teal!4,colframe=teal!35!black,title={Before you move on}]
An earthquake? (\textbf{\ja{じしん}}, \emph{jishin}.) To crumble? (\textbf{\ja{くずれる}}, \emph{kuzureru}.) To roll? (\textbf{\ja{ころがる}}, \emph{korogaru}.) Frightening? (\textbf{\ja{おそろしい}}, \emph{osoroshii}.) Evacuation? (\textbf{\ja{ひなん}}, \emph{hinan}.) A tsunami? (\textbf{\ja{つなみ}}, \emph{tsunami}.) To flow? (\textbf{\ja{ながれる}}, \emph{nagareru}.) To wash away? (\textbf{\ja{ながす}}, \emph{nagasu}.) To float? (\textbf{\ja{うく}}, \emph{uku}.) Once more? (\textbf{\ja{また}}, \emph{mata}.) Heavy rain? (\textbf{\ja{おおあめ}}, \emph{ōame}.) A flood? (\textbf{\ja{おおみず}}, \emph{ōmizu}.) To overflow? (\textbf{\ja{あふれる}}, \emph{afureru}.) A cliff? (\textbf{\ja{がけ}}, \emph{gake}.) Awful? (\textbf{\ja{たいへん}}, \emph{taihen}.) Heavy snow? (\textbf{\ja{おおゆき}}, \emph{ōyuki}.) To accumulate? (\textbf{\ja{つもる}}, \emph{tsumoru}.) A snowstorm? (\textbf{\ja{ふぶき}}, \emph{fubuki}.) A power cut? (\textbf{\ja{ていでん}}, \emph{teiden}.) To go numb with cold? (\textbf{\ja{こごえる}}, \emph{kogoeru}.) A flash of lightning? (\textbf{\ja{いなずま}}, \emph{inazuma}.) To strike? (\textbf{\ja{あたる}}, \emph{ataru}.) Violent? (\textbf{\ja{はげしい}}, \emph{hageshii}.) A tornado? (\textbf{\ja{たつまき}}, \emph{tatsumaki}.) To lean? (\textbf{\ja{かたむく}}, \emph{katamuku}.) A volcano? (\textbf{\ja{かざん}}, \emph{kazan}.) An eruption? (\textbf{\ja{ふんか}}, \emph{funka}.) A disaster? (\textbf{\ja{さいがい}}, \emph{saigai}.) To prevent? (\textbf{\ja{ふせぐ}}, \emph{fusegu}.) Big? (\textbf{\ja{おおきい}}, \emph{ōkii}.)
\end{tcolorbox}
